% =============================================================
%  MAT231-F27 Calculus II — Lecture 17
%  Unit 17: Implicit and inverse function theorems (statement
%  and use)
%
%  Build from the repository root:   npm run materials -- lecture-17
%  This writes lecture-17.pdf (slides, with overlays) and
%  lecture-17-handout.pdf (print version) into this folder.
% =============================================================
\documentclass{wosm-slides}

\usepackage{pgfplots}
\pgfplotsset{compat=1.18}

\newcommand{\R}{\mathbb{R}}
\newcommand{\norm}[1]{\lVert #1 \rVert}
\newcommand{\vx}{\mathbf{x}}
\newcommand{\vy}{\mathbf{y}}
\newcommand{\va}{\mathbf{a}}
\newcommand{\vb}{\mathbf{b}}
\newcommand{\vc}{\mathbf{c}}

\course{MAT231-F27}{Calculus II}
\title[L17 \textperiodcentered\ Implicit and inverse functions]{Implicit and inverse functions}
\subtitle{The implicit and inverse function theorems: statement and use}
\author{Dr.\ Ali Khudiyev}
\institute{WoSM Academy}
\date{Lecture 17\enspace\textperiodcentered\enspace Unit 17\enspace\textperiodcentered\enspace Fall 2027}

\begin{document}

\maketitle

% -------------------------------------------------------------
\begin{frame}{Last time}
  \begin{itemize}
    \item Lagrange: at an extremum of $f$ on $\{g = c\}$ with $\nabla g \ne \mathbf 0$, \
      $\nabla f = \lambda\,\nabla g$.
    \item The multiplier is the rate of change of the optimum: $M'(c) = \lambda$.
    \item Several constraints: $\nabla f = \sum_j \lambda_j \nabla g_j$, if the $\nabla g_j$ are
      linearly independent.
  \end{itemize}
  \vskip0.6em
  The proof borrowed one fact: if $\nabla g(\va) \ne \mathbf 0$, then near $\va$ the set
  $\{g = c\}$ is a smooth surface. Today: where that fact comes from, and what else it gives.
\end{frame}

\begin{frame}{This lecture}
  \tableofcontents
  \vskip1em
  {\small\color{wosm-muted}
  Reading: Marsden \& Tromba, \emph{Vector Calculus}, Section 3.5.\par}
  \vskip0.6em
  {\small This unit is about statements and their use. The proofs are longer than one
  lecture; we give the idea of each.\par}
\end{frame}

% =============================================================
\section{Implicit functions}

\begin{frame}{Curves given by equations}
  \begin{columns}[c]
    \begin{column}{0.56\textwidth}
      The circle $x^2 + y^2 = 1$ is not the graph of a function $y(x)$. But near most of its points
      it is:
      \begin{itemize}
        \item near $(0, 1)$: \ $y = \sqrt{1 - x^2}$;
        \item near $(0, -1)$: \ $y = -\sqrt{1 - x^2}$;
        \item near $(1, 0)$: \ \textbf{no} function $y(x)$: for $x$ just below $1$ there are two
          values of $y$, for $x > 1$ none. Solve for $x$ instead: $x = \sqrt{1 - y^2}$.
      \end{itemize}
      \pause
      With $F(x, y) = x^2 + y^2$: \ $F_y = 2y$ vanishes exactly at $(\pm 1, 0)$, where the tangent is
      vertical.
    \end{column}
    \begin{column}{0.42\textwidth}
      \centering
      \begin{tikzpicture}[scale=1.5]
        \draw[wosm-muted, very thin] (-1.35,0) -- (1.35,0) (0,-1.35) -- (0,1.35);
        \draw[wosm-text, thick] (0,0) circle (1);
        \draw[wosm-accent, very thick] ({cos(60)},{sin(60)}) arc (60:120:1);
        \draw[wosm-accent, densely dashed] (-0.5,0.75) rectangle (0.5,1.12);
        \draw[wosm-brick, very thick] ({cos(-30)},{sin(-30)}) arc (-30:30:1);
        \draw[wosm-brick, densely dashed] (0.8,-0.55) rectangle (1.15,0.55);
        \fill[wosm-accent] (0,1) circle (1.2pt);
        \fill[wosm-brick] (1,0) circle (1.2pt);
      \end{tikzpicture}
      \par\small Near $(0, 1)$ the circle is a graph over the $x$-axis; near $(1, 0)$ it is not.
    \end{column}
  \end{columns}
\end{frame}

\begin{frame}{Implicit differentiation, revisited}
  Suppose $F(x, y(x)) = c$ for a differentiable function $y(x)$. The chain rule gives
  \[ F_x + F_y\, y'(x) = 0, \qquad\text{so}\qquad y'(x) = -\frac{F_x}{F_y} \quad\text{where } F_y \ne 0 . \]
  \pause
  \begin{example}[The unit circle]
    $F = x^2 + y^2$: \ $y' = -\dfrac{2x}{2y} = -\dfrac xy$. \ At $\bigl(\frac35, \frac45\bigr)$ the slope is
    $-\frac34$. At $(\pm 1, 0)$ the formula breaks down --- exactly where $y(x)$ does not exist.
  \end{example}
  \pause
  Calculus~I used this formula assuming that $y(x)$ exists. The theorem says when it does.
\end{frame}

\begin{frame}{The implicit function theorem}
  \begin{theorem}[Implicit function theorem, two variables]
    Let $F$ be $C^1$ near $(a, b)$, with $F(a, b) = c$ and $F_y(a, b) \ne 0$. Then there are open
    intervals $I \ni a$, $J \ni b$ and a unique function $\varphi\colon I \to J$ such that, inside
    $I \times J$,
    \[ F(x, y) = c \quad\iff\quad y = \varphi(x) . \]
    $\varphi$ is $C^1$, with $\varphi'(x) = -F_x / F_y$ at $(x, \varphi(x))$; if $F$ is $C^k$, so is $\varphi$.
  \end{theorem}
  \pause
  {\small In words: where the level curve is not vertical ($F_y \ne 0$), it is locally the graph of a
  $C^1$ function of $x$.\par}
\end{frame}

\begin{frame}{Why it is true}
  Say $F_y(a, b) > 0$. By continuity, $F_y > 0$ on a small rectangle around $(a, b)$.
  \pause
  \begin{enumerate}
    \item For each fixed $x$, $y \mapsto F(x, y)$ is strictly increasing there. So for each $x$ there
      is \textbf{at most one} $y$ with $F(x, y) = c$.
    \pause
    \item $F(a, b - \delta) < c < F(a, b + \delta)$. By continuity in $x$, the same inequalities hold for
      all $x$ close to $a$.
    \pause
    \item By the intermediate value theorem, for each such $x$ there is \textbf{exactly one}
      $y = \varphi(x)$ between $b - \delta$ and $b + \delta$ with $F(x, \varphi(x)) = c$.
  \end{enumerate}
  \pause
  {\small That $\varphi$ is continuous and differentiable takes more work (Marsden \& Tromba, Section 3.5).
  Once we know $\varphi$ is differentiable, the formula for $\varphi'$ is the chain rule.\par}
\end{frame}

\begin{frame}{Example: the folium of Descartes}
  \begin{columns}[c]
    \begin{column}{0.56\textwidth}
      $F(x, y) = x^3 + y^3 - 6xy = 0$. \ $F_x = 3x^2 - 6y$, \ $F_y = 3y^2 - 6x$.
      \pause
      \begin{itemize}
        \item At $(3, 3)$: $F_x = F_y = 9$, so $y' = -1$. Tangent: $x + y = 6$.
        \pause
        \item $F_y = 0$ where $y^2 = 2x$: on the curve, at the origin and at
          $(2^{5/3}, 2^{4/3}) \approx (3.17, 2.52)$, a vertical tangent.
        \pause
        \item At $(0, 0)$, $F_x = F_y = 0$: the curve crosses itself, and is a graph in neither
          direction.
      \end{itemize}
      \pause
      {\small No formula $y(x)$ is needed to find slopes and tangent lines.\par}
    \end{column}
    \begin{column}{0.42\textwidth}
      \centering
      \begin{tikzpicture}[scale=0.62]
        \clip (-3.3,-3.3) rectangle (4.5,4.5);
        \draw[wosm-muted, very thin] (-3.3,0) -- (4.5,0) (0,-3.3) -- (0,4.5);
        \draw[wosm-rule, dashed] (-5.3,3.3) -- (3.3,-5.3);
        \draw[wosm-text, thick] plot coordinates {(0.000,0.000) (0.102,0.002) (0.203,0.007) (0.305,0.016) (0.407,0.028) (0.508,0.043) (0.610,0.062) (0.711,0.084) (0.812,0.110) (0.912,0.139) (1.012,0.172) (1.111,0.207) (1.210,0.246) (1.308,0.288) (1.405,0.333) (1.501,0.382) (1.595,0.433) (1.688,0.486) (1.780,0.543) (1.870,0.602) (1.958,0.664) (2.043,0.727) (2.127,0.793) (2.208,0.861) (2.287,0.930) (2.363,1.001) (2.436,1.073) (2.506,1.147) (2.572,1.221) (2.636,1.296) (2.696,1.371) (2.753,1.447) (2.806,1.522) (2.856,1.598) (2.902,1.672) (2.945,1.747) (2.983,1.820) (3.018,1.893) (3.050,1.964) (3.077,2.034) (3.101,2.103) (3.122,2.169) (3.139,2.234) (3.152,2.298) (3.163,2.359) (3.170,2.418) (3.174,2.474) (3.175,2.529) (3.173,2.581) (3.168,2.631) (3.161,2.679) (3.151,2.724) (3.139,2.767) (3.125,2.807) (3.108,2.845) (3.090,2.881) (3.070,2.914) (3.048,2.945) (3.025,2.973) (3.000,3.000) (2.981,3.018) (2.960,3.036) (2.939,3.052) (2.916,3.068) (2.892,3.083) (2.867,3.097) (2.840,3.110) (2.813,3.122) (2.784,3.133) (2.753,3.143) (2.722,3.152) (2.689,3.159) (2.655,3.165) (2.619,3.170) (2.582,3.173) (2.544,3.174) (2.505,3.175) (2.465,3.173) (2.423,3.170) (2.380,3.166) (2.336,3.159) (2.291,3.151) (2.245,3.141) (2.198,3.130) (2.150,3.116) (2.101,3.101) (2.051,3.083) (2.000,3.064) (1.948,3.043) (1.896,3.020) (1.843,2.995) (1.789,2.967) (1.735,2.938) (1.681,2.907) (1.626,2.874) (1.571,2.839) (1.516,2.802) (1.461,2.763) (1.405,2.723) (1.350,2.680) (1.295,2.636) (1.240,2.589) (1.186,2.541) (1.132,2.492) (1.078,2.440) (1.025,2.387) (0.973,2.333) (0.921,2.277) (0.870,2.219) (0.820,2.160) (0.771,2.100) (0.724,2.038) (0.677,1.976) (0.631,1.912) (0.587,1.847) (0.544,1.781) (0.502,1.714) (0.462,1.647) (0.423,1.578) (0.386,1.509) (0.350,1.439) (0.316,1.368) (0.283,1.297) (0.253,1.226) (0.223,1.154) (0.196,1.081) (0.170,1.008) (0.146,0.935) (0.124,0.862) (0.104,0.788) (0.085,0.714) (0.068,0.640) (0.053,0.566) (0.040,0.492) (0.029,0.418) (0.020,0.343) (0.012,0.269) (0.006,0.194) (0.002,0.120) (0,0)};
        \draw[wosm-text, thick] plot coordinates {(-3.210,1.531) (-3.067,1.414) (-2.929,1.304) (-2.797,1.200) (-2.668,1.103) (-2.544,1.011) (-2.424,0.925) (-2.307,0.843) (-2.192,0.767) (-2.080,0.695) (-1.971,0.627) (-1.864,0.563) (-1.758,0.503) (-1.654,0.447) (-1.552,0.395) (-1.450,0.346) (-1.350,0.301) (-1.251,0.259) (-1.153,0.220) (-1.055,0.184) (-0.958,0.152) (-0.861,0.123) (-0.765,0.097) (-0.669,0.074) (-0.573,0.055) (-0.477,0.038) (-0.382,0.024) (-0.286,0.014) (-0.191,0.006) (-0.095,0.002) (0.000,0.000)};
        \draw[wosm-text, thick] plot coordinates {(0.000,-0.001) (0.001,-0.064) (0.003,-0.127) (0.006,-0.190) (0.011,-0.253) (0.017,-0.316) (0.024,-0.379) (0.033,-0.442) (0.043,-0.505) (0.054,-0.568) (0.066,-0.632) (0.080,-0.695) (0.096,-0.759) (0.112,-0.822) (0.130,-0.886) (0.150,-0.950) (0.171,-1.014) (0.193,-1.078) (0.216,-1.143) (0.241,-1.208) (0.268,-1.273) (0.296,-1.339) (0.325,-1.405) (0.356,-1.471) (0.388,-1.538) (0.422,-1.606) (0.457,-1.674) (0.494,-1.743) (0.533,-1.812) (0.574,-1.882) (0.616,-1.953) (0.660,-2.025) (0.705,-2.098) (0.753,-2.172) (0.803,-2.246) (0.854,-2.323) (0.908,-2.400) (0.964,-2.479) (1.022,-2.559) (1.082,-2.641) (1.145,-2.725) (1.211,-2.810) (1.279,-2.897) (1.350,-2.987) (1.424,-3.079) (1.501,-3.173) (1.581,-3.270)};
        \draw[wosm-accent, thick] (1.5,4.5) -- (4.5,1.5);
        \fill[wosm-accent] (3,3) circle (2.5pt);
        \fill[wosm-brick] (3.175,2.520) circle (2.5pt);
        \fill[wosm-brick] (0,0) circle (2.5pt);
      \end{tikzpicture}
      \par\small $x^3 + y^3 = 6xy$, its tangent at $(3, 3)$, and the points where $F_y = 0$.
    \end{column}
  \end{columns}
\end{frame}

% =============================================================
\section{The implicit function theorem in general}

\begin{frame}{Surfaces given by an equation}
  For $F(x, y, z) = c$ with $F_z \ne 0$ at a point, the same theorem lets us solve for
  $z = \varphi(x, y)$ near it, and
  \[ \frac{\partial z}{\partial x} = -\frac{F_x}{F_z}, \qquad \frac{\partial z}{\partial y} = -\frac{F_y}{F_z} . \]
  \pause
  \begin{example}[The sphere $x^2 + y^2 + z^2 = 14$ at $(1, 2, 3)$]
    $F_z = 2z = 6 \ne 0$, so near $(1, 2, 3)$ the sphere is a graph $z = \varphi(x, y)$, with
    $z_x = -\frac{x}{z} = -\frac13$ and $z_y = -\frac yz = -\frac23$.
    Tangent plane: $z = 3 - \frac13 (x - 1) - \frac23 (y - 2)$, that is $x + 2y + 3z = 14$ --- the same
    as $\nabla F(\va) \cdot (\vx - \va) = 0$ (Lecture 11).
  \end{example}
\end{frame}

\begin{frame}{Systems of equations}
  $m$ equations in $n + m$ unknowns: \ $F(\vx, \vy) = \vc$, with $\vx \in \R^n$, $\vy \in \R^m$ and
  $F\colon \R^{n+m} \to \R^m$. We want to solve for $\vy$ in terms of $\vx$.
  \pause
  \begin{theorem}[Implicit function theorem]
    Let $F$ be $C^1$ near $(\va, \vb)$, with $F(\va, \vb) = \vc$, and suppose the $m \times m$ matrix
    $D_{\vy}F(\va, \vb) = \bigl(\partial F_i / \partial y_j\bigr)$ is invertible. Then near $(\va, \vb)$
    the solutions of $F(\vx, \vy) = \vc$ are exactly the points $(\vx, \varphi(\vx))$ of the graph of a
    $C^1$ function $\varphi$, and
    \[ D\varphi(\vx) = -\bigl(D_{\vy}F\bigr)^{-1} D_{\vx}F . \]
  \end{theorem}
  \pause
  {\small The formula: differentiate $F(\vx, \varphi(\vx)) = \vc$ by the chain rule (Lecture 12):
  $D_{\vx}F + D_{\vy}F\, D\varphi = 0$.\par}
\end{frame}

\begin{frame}{The linear model}
  A linear system $A\vx + B\vy = \vc$ with $B$ an invertible $m \times m$ matrix has the solution
  $\vy = B^{-1}(\vc - A\vx)$.
  \pause
  \begin{block}{The idea of the theorem}
    Near $(\va, \vb)$ a nonlinear system behaves like its linearisation ($A = D_{\vx}F$,
    $B = D_{\vy}F$). If that can be solved for $\vy$, so can the nonlinear system --- locally.
  \end{block}
  \pause
  \begin{itemize}
    \item $m$ independent equations in $n + m$ unknowns leave an $n$-dimensional solution set.
    \item One equation $g = c$ in $\R^n$ with $\nabla g(\va) \ne \mathbf 0$: some $g_{x_i}(\va) \ne 0$, so we
      can solve for $x_i$, and $\{g = c\}$ is a smooth hypersurface near $\va$ (the fact used in
      Lecture 16). Two equations in $\R^3$ with independent gradients: a smooth curve.
  \end{itemize}
\end{frame}

\begin{frame}{Example: a curve in space}
  The sphere $x^2 + y^2 + z^2 = 6$ meets the plane $x + y + z = 2$ in a curve through $(1, 2, -1)$.
  Can we solve for $y, z$ as functions of $x$?
  \pause
  \[ D_{(y, z)}F = \begin{pmatrix} 2y & 2z \\ 1 & 1 \end{pmatrix} = \begin{pmatrix} 4 & -2 \\ 1 & 1 \end{pmatrix},
     \qquad \det = 6 \ne 0 . \quad \text{Yes.} \]
  \pause
  Differentiate both equations in $x$:
  \[ 2x + 2y\, y' + 2z\, z' = 0, \qquad 1 + y' + z' = 0 . \]
  At $(1, 2, -1)$: \ $2 + 4y' - 2z' = 0$ and $y' + z' = -1$, so $y' = -\frac23$, $z' = -\frac13$.
  \pause
  \\[0.4em] Tangent direction $(1, -\frac23, -\frac13) \parallel (3, -2, -1)$, perpendicular to both
  gradients $(2, 4, -2)$ and $(1, 1, 1)$, as it should be.
\end{frame}

\begin{frame}{When the hypothesis fails}
  \begin{example}[The sphere $x^2 + y^2 + z^2 = 3$ and the plane $x + y + z = 3$]
    At $(1, 1, 1)$: \ $D_{(y, z)}F = \begin{pmatrix} 2 & 2 \\ 1 & 1 \end{pmatrix}$, \ $\det = 0$. The
    gradients $(2, 2, 2)$ and $(1, 1, 1)$ are parallel.
    \pause\\[0.3em]
    In fact the plane is at distance $3/\sqrt3 = \sqrt3$ from the origin: it \textbf{touches} the sphere,
    and the solution set is the single point $(1, 1, 1)$ --- not a curve.
  \end{example}
  \pause
  \begin{warning}[The theorem is one-directional]
    If $D_{\vy}F$ is singular, the theorem says nothing: the solutions may still be the graph of a
    function ($y^3 = x$ near the origin, where $F_y = 3y^2 = 0$), or not (above).
  \end{warning}
\end{frame}

% =============================================================
\section{The inverse function theorem}

\begin{frame}{Inverse functions in one variable}
  If $f$ is $C^1$ and $f'(a) \ne 0$, then near $a$, $f$ is strictly monotone, so it has an inverse,
  and
  \[ (f^{-1})'\bigl(f(a)\bigr) = \frac{1}{f'(a)} \qquad \text{(Calculus I, Unit 11).} \]
  \pause
  If $f'(a) = 0$, anything can happen:
  \begin{itemize}
    \item $x^2$ near $0$: not injective, so no inverse.
    \item $x^3$ near $0$: injective, but the inverse $\sqrt[3]{y}$ is not differentiable at $0$.
  \end{itemize}
  \pause
  In several variables, $f'(a)$ becomes the Jacobian matrix $DF(\va)$, and ``$\ne 0$'' becomes
  ``invertible'': $\det DF(\va) \ne 0$.
\end{frame}

\begin{frame}{The inverse function theorem}
  \begin{theorem}
    Let $F\colon U \to \R^n$ be $C^1$ on an open set $U \subseteq \R^n$, and let $DF(\va)$ be
    invertible. Then there are open sets $V \ni \va$ and $W \ni F(\va)$ such that $F\colon V \to W$ is
    a bijection, its inverse $G\colon W \to V$ is $C^1$, and
    \[ DG\bigl(F(\vx)\bigr) = DF(\vx)^{-1} \qquad (\vx \in V) . \]
  \end{theorem}
  \pause
  \begin{itemize}
    \item The formula is the chain rule applied to $G(F(\vx)) = \vx$: \ $DG\, DF = I$.
    \pause
    \item \textbf{Proof idea}: apply the implicit function theorem to $H(\vx, \vy) = F(\vx) - \vy = \mathbf 0$,
      solving for $\vx$ in terms of $\vy$. That needs $D_{\vx}H = DF(\va)$ to be invertible --- which
      is the hypothesis.
  \end{itemize}
\end{frame}

\begin{frame}{Example: polar coordinates}
  $F(r, \theta) = (r\cos\theta, \ r\sin\theta)$:
  \[ DF = \begin{pmatrix} \cos\theta & -r\sin\theta \\ \sin\theta & r\cos\theta \end{pmatrix}, \qquad
     \det DF = r . \]
  \pause
  \begin{itemize}
    \item $r \ne 0$: $F$ is locally invertible, and
      $DF^{-1} = \begin{pmatrix} \cos\theta & \sin\theta \\ -\frac{\sin\theta}{r} & \frac{\cos\theta}{r} \end{pmatrix}$.
      Its first row gives $\frac{\partial r}{\partial x} = \cos\theta = \frac xr$; the second,
      $\frac{\partial\theta}{\partial x} = -\frac{\sin\theta}{r} = -\frac{y}{x^2 + y^2}$ --- without a
      formula for $\theta(x, y)$.
    \pause
    \item $r = 0$: every $\theta$ is sent to the origin. No local inverse there.
    \pause
    \item \textbf{Local, not global}: $F(r, \theta + 2\pi) = F(r, \theta)$. Each point has a neighbourhood
      on which $F$ is injective, but $F$ is not injective on $\{r > 0\}$.
  \end{itemize}
  {\small $|\det DF| = r$ is the area factor of polar coordinates (Lecture 20).\par}
\end{frame}

\begin{frame}{Local invertibility everywhere is not enough}
  \begin{example}[$F(x, y) = (e^x\cos y, \ e^x\sin y)$]
    \[ \det DF = \det\begin{pmatrix} e^x\cos y & -e^x\sin y \\ e^x\sin y & e^x\cos y \end{pmatrix} = e^{2x} > 0
       \quad \text{everywhere,} \]
    so $F$ is locally invertible at every point. But $F(x, y + 2\pi) = F(x, y)$: not injective.
  \end{example}
  \pause
  \begin{itemize}
    \item In one variable, $f' \ne 0$ on an interval forces $f$ to be monotone, hence injective.
      In two variables there is no such argument.
    \item In complex notation this $F$ is $z \mapsto e^z$, and its local inverses are the branches
      of the logarithm.
  \end{itemize}
\end{frame}

\begin{frame}{Application: a robot arm}
  \begin{columns}[c]
    \begin{column}{0.58\textwidth}
      Two links of length $1$, joint angles $\theta_1, \theta_2$. The hand is at
      \[ F(\theta_1, \theta_2) = \begin{pmatrix} \cos\theta_1 + \cos(\theta_1 + \theta_2) \\
         \sin\theta_1 + \sin(\theta_1 + \theta_2) \end{pmatrix} . \]
      \pause
      A short computation gives $\det DF = \sin\theta_2$.
      \pause
      \begin{itemize}
        \item $\sin\theta_2 \ne 0$: the angles depend smoothly on the hand position; the hand can
          move in any direction.
        \item $\theta_2 = 0$ (stretched) or $\pi$ (folded): a \textbf{singular} configuration. From
          full stretch the hand cannot move straight outwards.
      \end{itemize}
    \end{column}
    \begin{column}{0.4\textwidth}
      \centering
      \begin{tikzpicture}[scale=1.15]
        \draw[wosm-rule, dashed] (0,0) circle (2);
        \draw[wosm-muted, very thin] (-0.3,0) -- (2.1,0);
        \draw[wosm-text, very thick] (0,0) -- ({cos(30)},{sin(30)}) -- ({cos(30) + cos(100)},{sin(30) + sin(100)});
        \fill[wosm-text] (0,0) circle (1.6pt);
        \fill[wosm-accent] ({cos(30)},{sin(30)}) circle (1.6pt);
        \fill[wosm-brick] ({cos(30) + cos(100)},{sin(30) + sin(100)}) circle (1.8pt);
        \draw[wosm-muted] (0.35,0) arc (0:30:0.35) node[right, font=\scriptsize, xshift=1pt, yshift=-2pt] {$\theta_1$};
        \draw[wosm-muted, densely dotted] ({cos(30)},{sin(30)}) -- ({1.45*cos(30)},{1.45*sin(30)});
        \draw[wosm-muted] ({cos(30) + 0.3*cos(30)},{sin(30) + 0.3*sin(30)}) arc (30:100:0.3)
          node[above right, font=\scriptsize, xshift=2pt] {$\theta_2$};
      \end{tikzpicture}
      \par\small The hand reaches the disc of radius $2$; its edge only with $\theta_2 = 0$.
    \end{column}
  \end{columns}
\end{frame}

\begin{frame}{Application: how a solution moves with a parameter}
  For each $p$, the equation $x^3 + x = p$ has exactly one real solution $x(p)$ ($x^3 + x$ is strictly
  increasing). There is a formula (Cardano's), but we do not need it.
  \pause
  \begin{itemize}
    \item $F(x, p) = x^3 + x - p$ has $F_x = 3x^2 + 1 > 0$, so $x(p)$ is $C^1$ (even $C^\infty$), with
      \[ x'(p) = -\frac{F_p}{F_x} = \frac{1}{3x^2 + 1} . \]
    \item At $p = 2$: $x = 1$, \ $x'(2) = \frac14$. \ So $x(2.1) \approx 1.025$ \ (true value $1.0245$).
  \end{itemize}
  \pause
  \begin{remark}
    The same idea tells how the minimiser of $f(\vx, p)$ moves with a parameter $p$: it solves
    $\nabla_{\vx} f = \mathbf 0$, and the implicit function theorem applies when the Hessian is
    invertible.
  \end{remark}
\end{frame}

% =============================================================
\section{Summary and practice}

\begin{frame}{Summary}
  \begin{itemize}
    \item \textbf{Implicit function theorem}: $F(\vx, \vy) = \vc$ can be solved locally as
      $\vy = \varphi(\vx)$ when $D_{\vy}F$ is invertible; then $D\varphi = -(D_{\vy}F)^{-1} D_{\vx}F$.
      For one equation: $y' = -F_x / F_y$.
    \item So level sets $\{g = c\}$ with $\nabla g \ne \mathbf 0$ are smooth hypersurfaces, and $m$
      equations with independent gradients define an $(n + m) - m = n$-dimensional set.
    \item \textbf{Inverse function theorem}: $\det DF(\va) \ne 0$ gives a local $C^1$ inverse, with
      $DG = DF^{-1}$.
    \item Both are \textbf{local}: polar coordinates and $e^z$ are locally but not globally invertible.
    \item Uses: tangent lines and planes without formulas; singular configurations; how solutions
      depend on parameters; change of variables (Lecture 20).
  \end{itemize}
\end{frame}

\begin{frame}{Practice}
  \small
  \begin{enumerate}
    \item At which points of the ellipse $x^2 - xy + y^2 = 3$ can we \emph{not} solve for $y$ as a function
      of $x$?
    \item $e^{xy} + x + y = 2$ defines $y(x)$ near $(0, 1)$. Find $y'(0)$.
    \item $xz + yz^3 + z = 3$ defines $z(x, y)$ near $(1, 1, 1)$. Find $z_x$ and $z_y$ there.
    \item Show that near $(1, 1, 1)$ the system $x^2 + y^2 + z^2 = 3$, $x + y + z = 3$ cannot be solved for
      $(y, z)$ by the theorem, and find all its solutions.
    \item Where is $F(u, v) = (u^2 - v^2, \ 2uv)$ locally invertible? Is it injective on
      $\R^2 \setminus \{\mathbf 0\}$?
    \item For polar coordinates, find $\partial r / \partial y$ and $\partial\theta / \partial y$ from $DF^{-1}$.
    \item Check that $\det DF = \sin\theta_2$ for the robot arm. Which hand positions are reached in
      singular configurations?
  \end{enumerate}
\end{frame}

\begin{frame}{Answers}
  \small
  \begin{enumerate}
    \item $F_y = 2y - x = 0$ and the equation give $3y^2 = 3$: \ the points $(2, 1)$ and $(-2, -1)$.
    \item $F_x = ye^{xy} + 1 = 2$, \ $F_y = xe^{xy} + 1 = 1$: \ $y'(0) = -2$.
    \item $F_x = z = 1$, \ $F_y = z^3 = 1$, \ $F_z = x + 3yz^2 + 1 = 5$: \ $z_x = z_y = -\frac15$.
    \item $\det\begin{pmatrix} 2y & 2z \\ 1 & 1 \end{pmatrix} = 2y - 2z = 0$ at $(1, 1, 1)$. The only solution is
      $(1, 1, 1)$: the plane touches the sphere.
    \item $\det DF = 4(u^2 + v^2)$: locally invertible except at the origin. Not injective:
      $F(-u, -v) = F(u, v)$. (It is $z \mapsto z^2$.)
    \item $\frac{\partial r}{\partial y} = \sin\theta = \frac yr$, \ $\frac{\partial\theta}{\partial y} = \frac{\cos\theta}{r} = \frac{x}{x^2 + y^2}$.
    \item Expand the $2 \times 2$ determinant and use $\sin(\alpha - \beta) = \sin\alpha\cos\beta - \cos\alpha\sin\beta$.
      $\theta_2 = 0$: the circle of radius $2$; \ $\theta_2 = \pi$: the origin.
  \end{enumerate}
\end{frame}

\begin{frame}{Next time}
  \textbf{Lecture 18 --- Double integrals} (Unit 18)
  \begin{itemize}
    \item Integrals over rectangles: Riemann sums and volume
    \item Fubini's theorem: computing by iterated integrals
    \item General regions, and changing the order of integration
  \end{itemize}
  \vskip1em
  {\small
  Reading: Strang \& Herman, \emph{Calculus, Volume 3}, Sections 5.1 and 5.2.\\[0.4em]
  Questions: the module channel on Discord, or office hours,
  Monday--Friday, 10:00--17:00 (Baku time, UTC+4),
  \href{mailto:ali.khudiyev@wosm.academy}{ali.khudiyev@wosm.academy}.\par}
\end{frame}

\end{document}
